% Proof fragment. Requires amsmath/amssymb, xurl and hyperref; no custom theorem environments.
% Prepared against ProveIt commit 17f92105e2bc4317feae6c219cb2143437b482ff.
% This is an ordinary mathematical proof, not a new kernel-checked endpoint.
\section[Uniform calibration for formula radices k >= 4]{Uniform calibration for formula radices $k\geq4$}
\label{sec:radix-calibration}

The formula multiplier can be changed from four to five or six without
changing the consistency strength of the numerical calculus. This is
not justified by replacing the numeral four in the old proof: the old
calculus identifies a formula's code with a fresh variable label, and
its propositional axiom has a numerical side condition. Both features
must be addressed explicitly. The construction below does so for every
fixed natural number $k\geq4$. The literal machines in the accompanying
packet instantiate it at $k=5$ and $k=6$.

\subsection[The numerical calculus at radix k]{The numerical calculus at radix $k$}

Let
\[
 p(a,b)=\frac{(a+b)(a+b+1)}2+b
\]
be Cantor pairing, with its usual primitive-recursive inverse. Define
\begin{align*}
 I_k(A,B)&=kp(A,B),&M_k(a,b)&=kp(a,b)+1,\\
 U_k(a,B)&=kp(B,a)+2,&E_k(a,b)&=kp(a,b)+3.
\end{align*}
The reversal in $U_k$ is deliberate. Code zero denotes truth; code one,
which equals $M_k(0,0)$, is the distinguished false formula. Put
\[
 N_k(B)=I_k(B,1),\qquad
 A\mathbin{\wedge_k}B=N_k(I_k(A,N_k(B))),\qquad
 \exists_k x_a B=N_k(U_k(a,N_k(B))).
\]
In formula displays, implication, negation, conjunction and existential
quantification refer to these numerical constructors.

The six retained heads are
\begin{align*}
 h_1(D,H)&=(\exists_k x_D\,U_k(D,H))\to H,\\
 h_2(D,A,H)&=U_k(D,U_k(A,H))\to U_k(A,U_k(D,H)),\\
 h_3(D,A,H)&=M_k(H,D)\to(E_k(H,A)\to M_k(A,D)),\\
 h_4(D,A,H)&=E_k(H,D)\to(M_k(A,H)\to M_k(A,D)),\\
 h_5(D,H)&=\exists_k x_D E_k(H,D),\\
 h_6(D,A,H)&=U_k(D,A)\to(U_k(D,A\to H)\to U_k(D,H)).
\end{align*}
The propositional head is
\[
 \rho_{8,k}(D,A,H,u)
   =((H\to u)\to D)\to((D\to H)\to(A\to H)),
 \qquad u\leq D.
\]
Define
\begin{align*}
 \Gamma_k(B)&=I_k(B,U_k(B,B)),\\
 V_k(a,b,c)&=M_k(a,b)\to U_k(c,M_k(a,b)),\\
 P_k(a,b,c)&=c\to M_k(c,b),\\
 S_k(a,b,c)&=U_k\bigl(b,
       (M_k(a,b)\to P_k(a,b,c))\wedge_k(M_k(c,b)\to c)\bigr),\\
 Q_k(a,b,c)&=U_k(c,S_k(a,b,c)\to N_k(P_k(a,b,c))),\\
 G_k(a,b,c)&=N_k(V_k(a,b,c)\to Q_k(a,b,c)).
\end{align*}
Let $\mathcal C_k$ have these six heads, the dominated propositional
heads, every $\Gamma_k(B)$, and $G_k(a,b,c)$ when $c\ne a,b$.
Its rules are modus ponens and generalization
$B\mapsto U_k(i,B)$ on derivations with no undischarged hypotheses.
The formulas themselves may have free variable labels. Failed guarded
emissions produce zero, which is already derivable since
$\rho_{8,k}(0,0,0,0)=0$.

\subsection{Hereditary recoding of formulas and variable labels}

Cantor pairing satisfies $a,b\leq p(a,b)$. Consequently, if
$n=kp(a,b)+r>0$, where $k\geq4$ and $0\leq r<k$, then $a,b<n$.
This gives the well-founded recursions used below.

Define $F_k(0)=0$. For $n>0$, write uniquely
$n=4p(a,b)+r$ with $0\leq r<4$, and set
\[
 F_k(n)=kp(F_k(a),F_k(b))+r.                 \tag{R1}
\]
This is course-of-values primitive recursion. For example, a primitive
recursion on $n$ stores the already computed sequence
$F_k(0),\ldots,F_k(n-1)$ and obtains the two required previous values
by primitive-recursive lookup. Thus $F_k$ is primitive recursive,
jointly in $k$ and $n$ on the stated parameter domain. PA proves its
defining identities by induction. It fixes $0,1,2,3$, and induction
on $n$ also proves $F_k(n)>0$ for $n>0$.

For injectivity, define a second primitive-recursive function $J_k$.
Set $J_k(0)=0$. For $m>0$, decode
$m=kp(a,b)+r$, $0\leq r<k$, and put
\[
 J_k(m)=
 \begin{cases}
  4p(J_k(a),J_k(b))+r,&r<4,\\
  0,&4\leq r<k.
 \end{cases}                                           \tag{R2}
\]
Whenever the recursive clause is used its arguments are smaller than
$m$. Induction on $n$, using positivity for $n>0$ and the uniqueness
of quotient, remainder and Cantor decoding, proves
\[
 J_k(F_k(n))=n.                                         \tag{R3}
\]
Therefore $F_k$ is injective. In particular, all numerical equality
and inequality guards among source labels are preserved. The function
$J_k$ is used only to establish injectivity; it is not asserted to
translate arbitrary $\mathcal C_k$ proofs back to radix four.

Directly from (R1), including the exceptional case $I_4(0,0)=0$, we get
\begin{align*}
 F_k(I_4(A,B))&=I_k(F_k(A),F_k(B)),\\
 F_k(M_4(a,b))&=M_k(F_k(a),F_k(b)),\\
 F_k(U_4(a,B))&=U_k(F_k(a),F_k(B)),\\
 F_k(E_4(a,b))&=E_k(F_k(a),F_k(b)).             \tag{R4}
\end{align*}
Thus the same map must be applied to a formula's numerical code and
to every numerical variable label occurring in it. Injectivity ensures
that binding, repeated labels and shadowing are respected: a binder
with old label $a$ binds exactly the occurrences that receive new
label $F_k(a)$. In particular,
\[
 F_k(\Gamma_4(B))
   =I_k(F_k(B),U_k(F_k(B),F_k(B)))
   =\Gamma_k(F_k(B)).                                  \tag{R5}
\]
This is the reason for using hereditary recoding. Recoding formula
constructors while holding the variable labels fixed would not
preserve the diagonal axiom. Equation (R4) also preserves each
retained head, each of $V,P,S,Q,G$, modus ponens and generalization;
injectivity preserves the guard on $G$.

\subsection{Removing the numerical-order obstruction}

The function $F_k$ need not preserve order. For instance,
\[
 40=4p(4,0)<4p(3,1)=44,
 \qquad F_5(40)=75>55=F_5(44).
\]
It would therefore be incorrect to translate a restricted
$\rho_{8,4}$ axiom into one restricted $\rho_{8,k}$ axiom merely
by recoding its four parameters.

There is, however, a fixed finite derivation of the unrestricted
schema from the restricted schema. The archived source
\texttt{M6/Logic/RestrictedImplication.lean} proves a theorem
\texttt{rho8All} for arbitrary objects $A$, an implication operation,
a relation $\preceq$, and proof constructors with the following
properties:
\begin{enumerate}
 \item modus ponens is available;
 \item the row $((c\to u)\to a)\to((a\to c)\to(b\to c))$
       is available if $u\preceq a$;
 \item $a\preceq a$;
 \item $u\preceq a$ implies $u\preceq(a\to b)$;
 \item $u\preceq b$ implies $u\preceq(a\to b)$.
\end{enumerate}
Its conclusion is that the displayed row has a derivation for
\emph{all} $a,b,c,u$, with no order hypothesis. Its intermediate
proofs derive weakening, Peirce's law, composition, exchange and the
usual implication distribution law. Every restricted leaf's side
condition is built solely by the last three listed operations.
The source is a finite proof template; it contains no search,
semantic completeness assumption or numerical-code substitution rule.

For $\mathcal C_k$, instantiate $\preceq$ by ordinary $\leq$ and
implication by $I_k$. The needed bounds follow from
\[
 a,b\leq p(a,b)\leq kp(a,b)=I_k(a,b).
\]
Transitivity gives the last two properties. This instantiation works
for every $k\geq4$ (indeed, the bound argument alone needs only
$k\geq1$). The fixed finite template is a primitive-recursive proof
constructor whose correctness PA verifies.

Translate a finite $\mathcal C_4$ derivation by (R4) at all ordinary
proof constructors. For each dominated axiom leaf, replace the leaf
by this unrestricted template at the recoded parameters. The
template's own restricted leaves all have valid numerical side
conditions at radix $k$. Induction on the original derivation gives
a primitive-recursive map $L_k$ with
\[
 \mathrm{PA}\vdash
 \operatorname{Prf}_{\mathcal C_4}(d,A)
 \ \longrightarrow\
 \operatorname{Prf}_{\mathcal C_k}(L_k(d),F_k(A)).        \tag{R6}
\]
In particular, $F_k(1)=1$ transfers every radix-four refutation to
a radix-$k$ refutation. No numerical monotonicity of $F_k$ is used.

\subsection{PA soundness, including the unused residues}

For $k>4$ there are natural numbers whose residue is not one of the
four constructor tags. They must receive an interpretation, since
arbitrary numerical command parameters can place them inside the
emitted formulas. Define a total syntactic translation $T_k$ into
arithmetic formulas as follows. Translate zero to the closed truth
$0=0$. For $n>0$, decode $n=kp(a,b)+r$:
\[
\begin{array}{c|l}
 r&T_k(n)\\ \hline
 0&T_k(a)\to T_k(b)\\
 1&\operatorname{Bit}(x_a,x_b)\\
 2&\forall x_b\,T_k(a)\\
 3&x_a=x_b\\
 4,\ldots,k-1&0=0.
\end{array}                                             \tag{R7}
\]
The auxiliary bound variables inside the bit predicate are renamed
above the finitely many source labels, or represented by de Bruijn
indices. Thus the labels $a,b$ in (R7) refer to the source variables,
with ordinary binding and shadowing semantics. The recursion is
primitive recursive because every recursive formula argument is
strictly smaller than $n$.

Every source label occurring in $T_k(B)$ is strictly below $B$.
This follows by induction from the component bound; zero and the
unused residues translate to closed truth. In particular $x_B$ is
not free in $T_k(B)$, so $T_k(\Gamma_k(B))$ has a uniform PA proof
by vacuity. The implication constructor has a uniform PA-provable
equivalence
\[
 T_k(I_k(A,B))\ \leftrightarrow\
                    (T_k(A)\to T_k(B)).                 \tag{R8}
\]
It is a syntactic identity except at $A=B=0$, where the proof is the
fixed tautology $(0=0)\leftrightarrow((0=0)\to(0=0))$.
This handles the numerical alias $I_k(0,0)=0$ explicitly.

Use the same bit interpretation as in the archived PA calibration.
PA proves
\[
 \operatorname{Bit}(s,z)\to z<s,
 \qquad\neg\operatorname{Bit}(z,z).
\]
For each arithmetic formula $\theta(z,\vec u)$ there is a
primitive-recursive supplier of a PA proof of finite comprehension:
\[
 \forall n\,\exists s<2^n\,\forall z<n\,
       (\operatorname{Bit}(s,z)\leftrightarrow\theta(z,\vec u)).
                                                               \tag{R9}
\]
It uses one PA induction instance on $n$. From a witness below
$2^n$, the next witness is either $s$ or $s+2^n$, according to
excluded middle on $\theta(n,\vec u)$. The proof supplier forms
that formula and induction instance syntactically; it never decides
$\theta$ and assumes no satisfaction or reflection principle.

For clarity, the guarded gate uses precisely the same finite
comprehension argument at every radix. Set $\phi=T_k(c)$ and assume
$c\ne a,b$. If the labels $a,b$ are distinct, use bound
$x_a+x_b+1$ and matrix
\[
  (\operatorname{Bit}(x_a,z)\vee z=x_b)\wedge\phi[z/x_b].
                                                               \tag{R10}
\]
If $a=b$, use bound $x_b+1$ and matrix
\[
  (z=x_b)\wedge\phi[z/x_b].                              \tag{R11}
\]
In the second case the first membership antecedent under the
quantifier for label $b$ is $\operatorname{Bit}(z,z)$ and is false.
In both cases the finite set witness makes the translation of
$S_k(a,b,c)$ true and satisfies the distinguished-point property
$P_k(a,b,c)$. The bound on set bits converts the bounded assertions
from (R9) into the required unbounded implications. Capture-avoiding
substitution preserves the original free values when choosing the
witness, and freshness of $c$ in $T_k(c)$ means the choice does not
alter $\phi$. We obtain a PA proof of $\neg T_k(Q_k(a,b,c))$.
The guard also makes $V_k$ a valid membership-vacuity instance, so
classical logic gives $T_k(G_k(a,b,c))$. When the guard fails the
emission is zero, with its fixed truth proof.

The other six heads have the same first-order proof suppliers as at
radix four. The propositional head is a classical tautology even
without its numerical restriction. Modus ponens uses (R8), and
generalization uses the ordinary PA rule. These arguments apply to
\emph{all} numerical parameter values, including unused-residue
formulas, since (R7) interprets those as truth. Finally,
\[
 T_k(1)=\operatorname{Bit}(x_0,x_0)
\]
has a fixed PA refutation, also uniformly under binders for label
zero. Hence the source negation and existential abbreviations have
their intended classical equivalences.

Recursing on the finite derivation therefore supplies a
primitive-recursive refutation translator $R_k$ such that
\[
 \mathrm{PA}\vdash
 \operatorname{Prf}_{\mathcal C_k}(d,1)
 \ \longrightarrow\
 \operatorname{Prf}_{\mathrm{PA}}(R_k(d),\ulcorner0=1\urcorner).
                                                               \tag{R12}
\]
This is a syntactic proof-checking theorem, so its universal
quantification includes the proof codes in nonstandard models of PA.
Combining (R6), (R12), and the archived PA calibration of
$\mathcal C_4$ gives
\[
 \mathrm{PA}\vdash
 \bigl(\exists d\,\operatorname{Prf}_{\mathcal C_k}(d,1)
       \ \leftrightarrow\
       \exists p\,\operatorname{Prf}_{\mathrm{PA}}
                              (p,\ulcorner0=1\urcorner)\bigr).
                                                               \tag{R13}
\]

\subsection{The ten-entry compiler works unchanged in form}

The numerical stack and command-word scheduler continue to use
ordinary Cantor pairing, without a multiplier. Only formula
construction uses $I_k,M_k,U_k,E_k$. Write $[a]w=p(a,w)$.
Shift tests $t+h=1$ and otherwise sends
$(t,w,d,h)$ to $(h,\operatorname{tl}(w),d,\operatorname{hd}(w))$.
For $v=d+x$, modus ponens compares $I_k(t,v)$ with $h$; on success
it returns $(v,[0]w,v,0)$ and on failure $(0,w,v,0)$.

The unchanged free-term emitter identity says that
$\mathsf{AX}(c,x,y)$, with $D=d+x,H=h+y,A=c$, pushes in order
\[
 t,\ U_k(d,h),\ \rho_{8,k}(D,A,H,d),\ h_1,h_2,h_3,h_4,h_5,h_6,
 \Gamma_k(A),
\]
leaves the guarded $G_k(H,D,A)$ on top, and clears the caches.
Let $q_1,\ldots,q_{10}$ denote these stack entries in reverse order.
The gate and $q_1,\ldots,q_8$ have residue zero modulo $k$;
$q_9=U_k(d,h)$ has residue two and is at least two. These facts
prove the same safe-shift identities as in the archived compiler:
\begin{align*}
 \mathsf{AX}(c,x,y)\mathsf S^9(t,w,d,h)
   &=(\rho_{8,k}(D,c,H,d),[t]w,0,U_k(d,h)),\\
 \mathsf{AX}(c,x,y)\mathsf S^{10}(t,w,d,h)
   &=(U_k(d,h),w,0,t).
\end{align*}
Every tested sum is different from one. No primality, parity, or
special identity of the number four is involved.

Consequently the same words give
\begin{align*}
 \mathrm{pop}&=\mathsf{AX}(0,0,0)\mathsf S^{11},\\
 \mathrm{push}&=\mathsf{AX}(0,0,0)\mathsf S^9\mathsf{MP}(0),\\
 \mathrm{discard}&=\mathsf{MP}(1)\mathsf{AX}(1,0,0)
                                      \mathsf S^{10}\mathsf{MP}(0).
\end{align*}
The MP comparison in push fails because implication and universal
codes have different residues. The two MP comparisons in discard
fail because $I_k(t,1)>0$ and $I_k(U_k(1,0),0)>0$.
For $i>0$, $\mathsf{MP}(i)$ at $(0,w,0,0)$ sets $d=i$ because
$I_k(0,i)>0$. These are the only numerical facts used by the
offset, selector, generalization and MP-suffix constructions.

For an actual restricted leaf with $u\leq D$, select its head by
\[
 \mathrm{discard}\,\operatorname{offset}(u)\,
 \mathsf{AX}(A,D-u,H)\mathsf S^9\mathrm{pop}.
\]
The subtraction here is performed on the \emph{target} numerical
parameters; $u+(D-u)=D$ follows from that leaf's valid side condition.
After (R6), the unrestricted-template expansion supplies such valid
restricted leaves, so all offsets are nonnegative.

The existing induction on a derivation now proves that its compiled
word either halts or returns the desired formula while preserving
the protected old stack and clearing the caches. A refutation gives
\[
 W_d=\mathsf{AX}(0,0,0)\,C_k(d)\,\mathsf S,
\]
which halts from every incoming command store. The leading AX clears
the two caches; it need not clear the carried stack. The alias
$p(0,0)=0$ and the possibility of an earlier halt are handled by
the same compiler identities.

There is an essential distinction here: $F_k$ embeds calculus
derivations, but it does not conjugate source stores or machine
trajectories. It preserves neither ordinary addition nor ordinary
Cantor-stack encoding. One must compile a \emph{new} target-radix
command word from the translated derivation. In particular, no
identity such as $F_k(D-u)=F_k(D)-F_k(u)$ is claimed or needed.

The command-word encoding and the scheduler are unchanged. Every
positive continuation pointer strictly decreases during its current
candidate, fresh candidates are $2,3,\ldots$, and an AX-prefixed
word has candidate code at least three. Thus PA proves that the
universally halting word is eventually reached, or a preceding
candidate has already halted. The arbitrary carried command store
causes no difficulty.

\subsection{Binding the changed arithmetic macro to the machines}

The fixed source block called \texttt{node61} is the four-increment
body in the formula-construction routine \texttt{Cons}. Its only
direct use is inside the decrementing loop \texttt{node62--node63}.
Replacing that body by $k$ increments changes the loop to
\[
 (T,Y)\longmapsto(T+k, Y-1)\quad(Y>0),
\]
with the zero branch exiting. After $i\leq Y_0$ iterations its
invariant is
\[
 T=T_0+ki,\qquad Y=Y_0-i,
\]
and all unrelated registers are preserved. Induction on $Y_0$
proves termination with $(T_0+kY_0,0)$. Counting the decrement,
$k$ increments and return jump gives exactly
$(k+2)Y_0+1$ primitive source steps, including the final zero test.

The preceding unpair and Cantor-pair macros are unchanged. For the
following source-macro contract, the four displayed coordinates are
$(T,W,X,Y)$, not the command-store coordinates $(t,w,d,h)$ used above.
The scratch registers satisfy $X=Y=0$ on entry. The new Cons routine sends
\[
 (t,[a]w,0,0)\longmapsto(I_k(a,t),w,0,0),
\]
preserving every other register.
The fixed additional increments for tags one, two and three remain
unchanged, so the same free-term emitter calculation produces the
heads displayed above, now with the radix-$k$ constructors. All
formula-construction uses, including the MP comparison, pass through
this Cons routine. Ordinary stack push and the scheduler's Cantor
coding do not pass through the formula-multiplier loop.

The source-binding check makes that separation explicit: there are
exactly $61$ occurrences of Cons and exactly $61$ occurrences of its
multiplier body, at the corresponding fixed relative path in each
Cons occurrence. There are no other uses of that body. Each Cons
has the same $21$-site unpair/pair prefix followed by one decrement,
$k$ increments and one back-edge jump. All $3429$ source sites
outside the old increment bodies retain their operations and every
successor edge under the checked ordinal relocation. Thus the source
has $3673+61(k-4)$ sites: $3734$ at radix five and $3795$ at radix six.

This loop induction is the changed arithmetic-macro proof. The
archived file \nolinkurl{M6/Source/AXTimes4.lean} states its radix-four
instance, with cost $6Y_0+1$, and
\nolinkurl{M6/Source/AXCons.lean} composes it with unpair and pairing.
Their fixed multiplier lemma cannot simply be cited as already
proving the new macro; the preceding induction supplies the required
generalization. The candidate's source expansion and literal
primitive contracts must additionally bind the new $k$-increment
body and its successors to the exact table.

For $k=5,6$, the supplied finite certificates do this final binding,
including the positive duration of every source primitive and the
complete blank-tape startup frame. PA formalizes their finite row
equalities and the register-loop inductions. The same sound-store
annotation used in the archived calibration works with $T_k$:
both formula registers and every canonical stack entry carry PA
proofs, successful MP combines them, AX uses the suppliers above,
and a halting Shift exposes code one in one of the two registers.
The zero-stack alias is handled by its fixed truth annotation.

Under $\operatorname{Con}(\mathrm{PA})$ this invariant prevents a
command halt. Positive physical simulation can therefore be
iterated past any proposed finite halting time, contradicting
determinism. Conversely, a PA refutation is transformed by (R13),
compiled into $W_d$, reached by the scheduler, and converted by the
physical contracts into a finite blank-tape halt. These operations
are primitive recursive and their correctness is provable in PA.
For the delivered exact table $M_k$, $k=5,6$, this gives
\[
 \boxed{\quad
 \mathrm{PA}\vdash
   \bigl(\neg\operatorname{Halt}_{\mathrm{blank}}(M_k)
          \leftrightarrow\operatorname{Con}(\mathrm{PA})\bigr).
 \quad}                                                \tag{R14}
\]
The statement is an ordinary mathematical PA-provability result
with executable finite certificates for the candidate-specific
program facts. The new radix endpoints are not claimed to have
been exported as newly kernel-checked Lean theorems.